\section*{Проверка статистических ошибок и качества аппроксимации}
Для отношения фиксированной аппроксимации $F$ к измеренной корреляционной
функции $C$ распространение ошибки знаменателя даёт
\[
    r=F/C,\qquad \sigma_r=\frac{|F|}{C^2}\sigma_C.
\]
Например, $F=1.5$, $C=1.2$, $\sigma_C=0.06$ дают $r=1.25$ и
$\sigma_r=0.0625$. Гистограмма вычисленной кривой имеет нулевую ошибку
при фиксированных параметрах и весах проекции. Это не утверждение
об отсутствии неопределённости оценённых параметров: её учёт требует
ковариационной матрицы и связи модели с теми же измерениями.
При $C=0$ отношение не определено.

\subsection*{Влияние округления}
В контрольном примере $n=10000$, $A=10000.0049$,
$Q_A=A^2/n+0.01$ используется
$\sigma_C^2=(Q_A-A^2/n)/(n(n-1))$.
Округление только $A$ до двух знаков даёт $C=1$ и
$\sigma_C=1.40720\cdot10^{-5}$ вместо $C=1.00000049$ и
$\sigma_C=1.00005\cdot10^{-5}$.
Изменение среднего на $4.9\cdot10^{-7}$ увеличивает ошибку на $40.7\%$.
Для трёх знаков ошибка равна $0.99000\cdot10^{-5}$; четыре и шесть
знаков сохраняют исходный результат. Это проверка десятичного округления,
а не воспроизведение накопления в гистограмме одинарной точности.
Округление для представления результатов выполняется после расчётов.
Дополнительно выполнены 10000 заполнений одного бина ROOT чередующимися
весами $1.00100049$ и $0.99900049$ при включённом \texttt{Sumw2}.
\texttt{TH1F} сохраняет сумму 10000, а \texttt{TH1D} --- 10000.0049.
Сырые ошибки среднего равны соответственно $1.40720\cdot10^{-5}$
и $1.00005\cdot10^{-5}$, хотя \texttt{Sumw2} хранится в двойной точности
в обоих случаях. Это показывает механизм потери точности при накоплении,
который действует и для отдельных ячеек трёхмерных гистограмм.
Сырые значения здесь вычислены до проверки разрешимости дисперсии:
в программе значение в пределах границы округления получает нулевую
записанную ошибку и исключается из аппроксимации.

\subsection*{Ручной расчёт качества аппроксимации}
Для пяти значений $C_i=(1.0,1.1,1.2,1.3,1.4)$ с ошибками $0.1$
подгоняется одна константа $F=1.2$. Вклады
$(C_i-F)^2/\sigma_i^2$ равны $(4,1,0,1,4)$, поэтому
\[
    \chi^2=10,\qquad \mathrm{ndf}=5-1=4,\qquad
    \chi^2/\mathrm{ndf}=2.5.
\]
Для трёхмерной КФ число степеней свободы определяется числом
использованных ячеек минус число свободных параметров.
Замороженные параметры не вычитаются. При обычном взвешенном
методе наименьших квадратов ячейки с нулевой ошибкой исключаются.
При интегральной аппроксимации сравнивается среднее модели внутри ячейки.
Приведённый пример проверяет алгоритм и не является результатом
для данных столкновений.

\subsection*{Проверка опции ROOT B}
Документация \texttt{TH1::Divide} определяет \texttt{B} как расчёт
биномиальных ошибок. Для числа прошедших отбор $k$ из $n$ испытаний
это соответствует $\sigma_{k/n}^2=p(1-p)/n$.
Для текущего среднего веса $C=\sum w_i/n$ числитель и знаменатель
строятся по общим парам, однако веса не обязаны быть индикаторами
прохождения отбора. При независимых весах оценка дисперсии среднего
равна $(\sum w_i^2-(\sum w_i)^2/n)/(n(n-1))$.
Контрольные результаты для 100 пар:
\begin{center}
\begin{tabular}{l|c|c|c}
Веса & $C$ & $\sigma$ ROOT B & $\sigma$ среднего\\ \hline
50 нулей и 50 единиц & 0.5 & 0.05 & 0.0502519\\
50 нулей и 50 двоек & 1 & 0 & 0.100504\\
Все веса равны 2 & 2 & 0.282843 & 0
\end{tabular}
\end{center}
Первая строка является биномиальной моделью; отличие оценки среднего
обусловлено поправкой $n/(n-1)$. Остальные строки демонстрируют,
что опция \texttt{B} не заменяет оценку ошибки произвольного среднего веса.
Для эффективности отбора \texttt{B} применима; для текущего среднего
веса её применение не обосновано. Событийные зависимости пар требуют
отдельного учёта ковариаций или повторной выборки по событиям.
Источник: \url{https://root.cern.ch/doc/v634/classTH1.html}.
Расчёты воспроизводятся макросом \texttt{tools/statistics\_review.C}.
